\begin{casebox}
\textbf{知识点小案例：随机分支为什么会把热循环拖慢。}
先看一个二值分类循环：统计数组中大于阈值的元素。若输入全小于阈值，分支方向稳定；若输入按 0、1、0、1 交替，预测器仍可能学到模式；若输入由随机数生成，分支方向接近不可预测。三组输入做同一件事，结果完全一样，但流水线看到的是三种前端和控制流形状。由特例推出规律：分支预测不是“if 快不快”的抽象问题，而是输入分布能否给预测器稳定历史。工程报告必须同时保存输入分布、checksum、分支版和 branchless 版反汇编；没有这些证据，不能把波动简单归因于“CPU 慢”。
\end{casebox}
